\begin{Exercise}[label={ex:operazioni}]
    Scrivere un programma che legga una sequenza di numeri interi positivi, terminata dal valore 0.
    Il programma deve calcolare, il numero totale di valori inseriti (escluso lo 0), la somma di tutti i numeri,
    il massimo e il minimo valore inserito e la media dei valori inseriti.

    Scrivere le funzioni:
    \begin{itemize}
        \item leggiNumero(): che legge e restituisce un numero naturale
        \item massimo(): restituisce il maggiore tra due numeri
        \item minimo(): restituisce il minore tra due numeri
        \item calcolaMedia(): restituisce la media dei numeri letti
    \end{itemize}
\end{Exercise}

\begin{Answer}[ref={ex:operazioni}]
  	\lstinputlisting[style=cstylesmall]{src/operazioni.c}
\end{Answer}